
\section{Surfaces including k-gons}

In this section, we will consider $k$-gons in $\mathbb{P}^3$ and their relationship with hypersurfaces.

\begin{definition}
    A $k$-gon is a union of $k$ lines $l_1, l_2, \dots, l_k$ such that $l_i$ intersects $l_{i+1}$ for $1 \leq i < k$, and $l_k$ intersects $l_1$. We will write $C_k = \bigcup_{i=1}^k l_i$ and denote the intersection points as $p_i = l_i \cap l_{i+1}$ for $1 \leq i < k$, and $p_k = l_k \cap l_1$.
\end{definition}

\begin{proposition}
    Let $C_k = \bigcup_{i=1}^k l_i$ be a $k$-gon in $\mathbb{P}^3$. Then we can find a hypersurface of degree $d$ such that $C_k \subseteq S_d$ when $k \leq \frac{(d+3)(d+2)(d+1)}{6d}$.
\end{proposition}

\begin{proof}

    We can look at the zero cohomology of the ideal sheaf $\mathcal{I}_{C_k}$ of the $k$-gon to see if the $k$-gon can be contained in a hypersurface. We will denote the union of the lines as $C_k = \bigcup_{i=1}^k l_i$. 

    We have the short exact sequence
    \[0 \to \mathcal{I}_{C_k} \to \mathcal{O}_{\mathbb{P}^3}\to \mathcal{O}_{C_k} \to 0.\]

    We twist this by $\mathcal{O}_{\mathbb{P}^3}(d)$ to get
    \[0 \to \mathcal{I}_{C_k}(d) \to \mathcal{O}_{\mathbb{P}^3}(d) \to \mathcal{O}_{C_k}(d) \to 0.\]

    From this we get the long exact sequence in cohomology:

    \begin{tikzcd}
    0 \rar & H^0(\mathbb{P}^3, \mathcal{I}_{C_k}(d)) \rar & H^0(\mathbb{P}^3, \mathcal{O}_{\mathbb{P}^3}(d)) \rar &  H^0(C_k, \mathcal{O}_{C_k}(d)) \ar[out=-20, in=160]{dll} \\
    & H^1(\mathbb{P}^3, \mathcal{I}_{C_k}(d)) \rar & H^1(\mathbb{P}^3, \mathcal{O}_{\mathbb{P}^3}(d)) \rar & H^1(C_k, \mathcal{O}_{C_k}(d)) \ar[out=-20, in=160]{dll} \\
    & H^2(\mathbb{P}^3, \mathcal{I}_{C_k}(d)) \rar & H^2(\mathbb{P}^3, \mathcal{O}_{\mathbb{P}^3}(d)) \rar & H^2(C_k, \mathcal{O}_{C_k}(d)) \ar[out=-20, in=160]{dll}\\
    & H^3(\mathbb{P}^3, \mathcal{I}_{C_k}(d)) \rar & H^3(\mathbb{P}^3, \mathcal{O}_{\mathbb{P}^3}(d)) \rar & H^3(C_k, \mathcal{O}_{C_k}(d)) \rar & 0
    \end{tikzcd}

    We know that $H^i(\mathbb{P}^3, \mathcal{O}_{\mathbb{P}^3}) = 0$ for $i > 0$. This means we can focus on the first part of the sequence:

    \[0 \to H^0(\mathbb{P}^3, \mathcal{I}_{C_k}) \to H^0(\mathbb{P}^3, \mathcal{O}_{\mathbb{P}^3}) \to H^0(C_k, \mathcal{O}_{C_k}) \to H^1(\mathbb{P}^3, \mathcal{I}_{C_k}) \to 0.\] 


    Since $h^0$ is additive, we get

    \[
    h^0(\mathbb{P}^3, \mathcal{I}_{C_k}) = h^0(\mathbb{P}^3, \mathcal{O}_{\mathbb{P}^3}) + h^1(\mathbb{P}^3, \mathcal{I}_{C_k}) - h^0(C_k, \mathcal{O}_{C_k}). 
    \]

    By Thm 18.4 in (EO) we get that $h^0(\mathbb{P}^3, \mathcal{O}_{\mathbb{P}^3}) = \binom{d+3}{3}$. $h^0(C_k, \mathcal{O}_{C_k}(d))$ can be computed by considering the short exact sequence: 
    \[
    0 \to \mathcal{O}_{C_k}(d) \to \bigoplus_{i=1}^k \mathcal{O}_{l_i}(d) \to \bigoplus_{i=1}^n \mathcal{O}_{p_i} \to 0.
    \]

    This gives us $h^0(C_k, \mathcal{O}_{C_k}) = k \cdot \binom{d+1}{1} - k = kd$.

    \[
    h^0(\mathbb{P}^3, \mathcal{I}_{C_k}) \geq \binom{d+3}{3} - kd.
    \]

    Meaning we can guarantee a degree $d$ surface when $\binom{d+3}{3} - kd \geq 0$, or equivalently when $k \leq \frac{(d+3)(d+2)(d+1)}{6d}$.

\end{proof}

To get an exact value for $h^0(\mathbb{P}^3, \mathcal{I}_{C_k})$, we would need to compute $h^1(\mathbb{P}^3, \mathcal{I}_{C_k})$. This issue will be explored in section ???. 


\begin{remark}
    Note that this argument does not guarantee that the surface is smooth, only that it exists.
\end{remark}

\section{Computational results}

We have done computational experiments in Macaulay2 to get a better understading of the inequality above. We have tested the existence and smoothness of hypersurfaces including $k$-gons for various degrees $d$ and number of lines $k$. The code used for these experiments is included in the following snippet. We will discuss the details of the code below, along with its results.

\lstinputlisting[caption={k-gon}, label={lst:k-gon}]{../Progging/ngon.m2}

We'll go through the code step by step. 

Step 1: We begin by defining the field and the polynomial ring in four variables over this field. We then define the degree $d$ of the surfaces we want to consider, and the number of points to use, which will be the same as the number of lines in the k-gon. 

Step 2: We define a function to pick a given amount of random points in our quotient ring, ensuring they are not zero. We then apply the function to get the points, and gather the points into pairs of subsequent points so that we can later define lines between these.

Step 3: Next, we define a function for turning two points into a matrix. The function takes two points and returns a matrix, as follows:

\[((a,b,c,d), (d,e,f,g)) \mapsto 
\begin{pmatrix}
    x_0 & x_1 & x_2 & x_3 \\
    a & b & c & d \\
    d & e & f & g 
\end{pmatrix}\]

Step 4: Now we can define the ideal of the line between the two points by taking the $3 \times 3$ minors of this matrix. 

Step 5: Next we take the intersection of all the line ideals, to get the ideal of $Z = \bigcup l_i$, 

Step 6: We choose a random member of $I_Z(d)$ and compute if this defines an empty, singular, or smooth surface.

Results: We have been interested in finding the maximum number of lines $k$ such that we can find a smooth surface of degree $d$ including a $k$-gon. The results are summarized in the following table:


\section{Cohomology of a 4-gon}

As mentioned earlier, to get an exact value for $h^0(\mathbb{P}^3, \mathcal{I}_{C_k})$, we would need to compute $h^1(\mathbb{P}^3, \mathcal{I}_{C_k})$. We will now explore this issue for the case of a 4-gon and looking at a suitable resolution of the ideal sheaf of the 4-gon. First we claim that the intersection points of a 4-gon can, by an automorphism of $\mathbb{P}^3$, be expressed as the points $(1:0:0:0), (0:1:0:0), (0:0:1:0), (0:0:0:1)$. Then the 4-gon can be expressed as the union of the lines $l_1 = V(x_2, x_3), l_2 = V(x_0, x_3), l_3 = V(x_0, x_1), l_4 = V(x_1, x_2)$. Now the ideal of the 4-gon is given by $I_{C_4} = (x_2, x_3) \cap (x_0, x_3) \cap (x_0, x_1) \cap (x_1, x_2) = (x_1x_3, x_0x_2)$. We justify our chose of points in the following lemma.

\begin{lemma}
    Let $C_4$ be a 4-gon in $\mathbb{P}^3$. Then there exists an automorphism of $\mathbb{P}^3$ such that the intersection points of $C_4$ are given by the points $(1:0:0:0), (0:1:0:0), (0:0:1:0), (0:0:0:1)$.
\end{lemma}

\begin{proof}
    Let $p_1, p_2, p_3, p_4$ be the intersection points of $C_4$. To be in general position in $\mathbb{P}^3$, we require that no three of the points are collinear, and that the points do not all lie in a plane. We can then find a projective transformation sending $p_1$ to $(1:0:0:0)$, $p_2$ to $(0:1:0:0)$, and $p_3$ to $(0:0:1:0)$. Since no three of the points are collinear, $p_4$ cannot be sent to a point on the line between any two of the other points. Furthermore, since the points do not all lie in a plane, $p_4$ cannot be sent to a point in the plane spanned by any three of the other points. Thus, we can send $p_4$ to $(0:0:0:1)$.
\end{proof}


We will now look at a resolution of the ideal sheaf of the 4-gon. We have that $I_{C_4} = (x_1x_3, x_0x_2)$. This ideal has the following free resolution:

\[0 \to \mathcal{O}_{\mathbb{P}^3}(-4) \xrightarrow{\begin{pmatrix}
    -x_0x_2 \\
    x_1x_3
\end{pmatrix}} \mathcal{O}_{\mathbb{P}^3}(-2)^{\oplus 2} \xrightarrow{\begin{pmatrix}
    x_1x_3 & x_0x_2
\end{pmatrix}} I_{C_4} \to 0.\]